\documentclass[11pt]{article}
\usepackage[a4paper,margin=26mm]{geometry}
\usepackage{amsmath,amssymb,amsthm,booktabs,array}
\usepackage[T1]{fontenc}
\usepackage[colorlinks=true,urlcolor=blue,citecolor=blue,linkcolor=blue]{hyperref}
\newtheorem{theorem}{Theorem}
\newtheorem{corollary}[theorem]{Corollary}
\newcommand{\aq}{a_q}
\title{Barker and Hardin's cyclic-repetition formulas\\and an extension of Fried's enumeration}
\author{Haobo Ma\\\small ChronoAI Pte Ltd, Singapore
\and Wenlin Zhang\\\small National University of Singapore\\\small ChronoAI Pte Ltd, Singapore}
\date{14 November 2026}
\begin{document}
\maketitle
\begin{abstract}
Extending Fried's repetition-word enumeration (arXiv:2410.07237v3,
26 October 2024, Theorem~10), we prove
Colin Barker's two conjectured formulas in the unpaginated FORMULA
section of OEIS A268943, revision 8 (17 January 2019), for the arrays introduced
by R. H. Hardin. The labels of successive adjacent equal pairs must advance
cyclically through an eight-letter alphabet. A three-state decomposition proves
the rational generating function and the radical expression for every positive
length. The same decomposition works for every alphabet size at least two,
establishing the empirical column recurrences and the first seven row
polynomials in Hardin's related table A268944. More generally, the enumeration
depends only on the common indegree of a loopless relation allowed between
successive event labels; Fried's unequal-label rule is one case. We also
count words by their number of equality events.
\end{abstract}

\section{Words and their equality events}
Fix an integer $q\geq2$ and write $\mathcal A_q=\{0,1,\ldots,q-1\}$.
For a word $w=(w_1,\ldots,w_n)$ over $\mathcal A_q$, an equality event is an
index $i<n$ with $w_i=w_{i+1}$; its label is $w_i$. List these labels in
increasing order of $i$, retaining overlapping pairs. The word is admissible
if each label after the first is the preceding label plus one modulo $q$.
An empty or singleton list of labels satisfies this condition. Thus $000$
is not admissible, whereas $(6,4,6,5,4)$ is admissible over eight letters.
Only the labels advance cyclically: the word itself is linear, and words
are not identified under rotation or relabeling.

Let $\aq(n)$ count admissible words of length $n$, with $\aq(0)=1$ for the
empty word, and put $S_q(x)=\sum_{n\geq1}\aq(n)x^n$. This zero-constant-term
convention matches the offset $1,1$ of A268943~\cite{hardin8}.
Its empirical recurrence and Barker's two conjectures concern $q=8$;
Hardin's table A268944 uses $T(n,k)=a_{k+1}(n)$ for $n,k\geq1$~\cite{hardintable}.
The FORMULA block in A268943 is headed ``Conjectures from Colin Barker,
Jan 17 2019''; Corollary~\ref{cor:barker} proves its two displayed expressions.
Fried~\cite[Section 3.2, Theorem 4]{fried}, first in
\cite[Section 2.8, Theorem 10]{friedpre}, counts words whose successive
event labels are unequal, rather than cyclic successors. The two restrictions
agree for $q=2$. For $q\geq3$ they differ: $(0,0,2,2)$ satisfies Fried's
condition but not ours. We give the cyclic-successor enumeration uniformly
in $q$, and recover the known binary formulas as its two-letter case.

\begin{theorem}\label{thm:general}
For every integer $q\geq2$,
\begin{equation}\label{eq:gf}
 S_q(x)=\frac{qx\bigl(1-(q-3)x-(q-1)x^2\bigr)}
 {(1-(q-1)x)(1-(q-2)x-qx^2)}.
\end{equation}
The initial values are $\aq(0)=1$, $\aq(1)=q$, $\aq(2)=q^2$, and
\begin{equation}\label{eq:rec}
 \aq(n+3)=(2q-3)\aq(n+2)+(-q^2+4q-2)\aq(n+1)
                -q(q-1)\aq(n)\qquad(n\geq0).
\end{equation}
\end{theorem}

\begin{proof}
For nonempty admissible words, use three disjoint states: $U$ has no equality
event; $A$ has first event label equal to the successor modulo $q$ of the
first letter; $B$ has an event but does not satisfy the condition for $A$.
Let $U_n,A_n,B_n$ be their respective counts at length $n\geq1$.
They start at $(U_1,A_1,B_1)=(q,0,0)$.

Prepend a letter $c$ to a word starting with $b$.
From $U$, the $q-1$ choices $c\ne b$ remain in $U$, and $c=b$ creates
the first event, of label $b$, putting the word in $B$.
Now suppose the first existing event has label $r$.
If $c=b$, a new event of label $b$ precedes $r$; this is permitted exactly
when $r=b+1\pmod q$, that is, for a word in $A$, and the resulting word
belongs to $B$.
If $c\ne b$, no event is added and admissibility is unchanged. The result
belongs to $A$ exactly when $c=r-1\pmod q$. For a word in $A$ this choice
is $b$ and is excluded, so all $q-1$ unequal choices belong to $B$.
For a word in $B$ there is exactly one allowed choice producing $A$ and
$q-2$ producing $B$; the equal choice is forbidden.
Removing the first letter deletes at most the first event, so an admissible
word has an admissible suffix. The first letter and suffix are unique;
therefore these cases give the complete transition
\begin{equation}\label{eq:states}
 U_{n+1}=(q-1)U_n,\qquad A_{n+1}=B_n,\qquad
 B_{n+1}=U_n+qA_n+(q-2)B_n.
\end{equation}

In the formal power series ring $\mathbb Q[[x]]$, set
$U(x)=\sum_{n\geq1}U_nx^n$ and define $A(x),B(x)$ similarly.
Equation~\eqref{eq:states} gives
\[
 U(x)=\frac{qx}{1-(q-1)x},\qquad A(x)=xB(x),\qquad
 B(x)=xU(x)+qxA(x)+(q-2)xB(x).
\]
Hence $B(x)=xU(x)/(1-(q-2)x-qx^2)$, and
$S_q(x)=U(x)+(1+x)B(x)$ gives~\eqref{eq:gf}. All denominators have
constant coefficient one, so these divisions are well-defined formally.
Adding the empty word yields
\begin{equation}\label{eq:fullgf}
 1+S_q(x)=\frac{1+(3-q)x+(2-q)x^2}
 {1-(2q-3)x+(q^2-4q+2)x^2+q(q-1)x^3}.
\end{equation}
Coefficient comparison at degrees zero, one and two gives the stated
initial values, and comparison at degree $n+3$ gives~\eqref{eq:rec}.
\end{proof}

\section{Barker's eight-letter formulas}
\begin{corollary}\label{cor:barker}
For $n\geq1$, the count $a_8(n)$ in A268943 has generating function
\[
 S_8(x)=\frac{8x(1-5x-7x^2)}{(1-7x)(1-6x-8x^2)}
\]
and satisfies
\begin{equation}\label{eq:radical}
 a_8(n)=\frac{-272\,7^n+(153-37\sqrt{17})(3-\sqrt{17})^n
                    +(153+37\sqrt{17})(3+\sqrt{17})^n}{34}.
\end{equation}
In particular, $a_8(n+3)=13a_8(n+2)-34a_8(n+1)-56a_8(n)$ for $n\geq0$.
\end{corollary}
\begin{proof}
The generating function and recurrence are Theorem~\ref{thm:general} at
$q=8$. The characteristic polynomial of the recurrence factors as
\[
 t^3-13t^2+34t+56=(t-7)(t^2-6t-8).
\]
Its roots are $7,3-\sqrt{17},3+\sqrt{17}$. Each corresponding power sequence,
and therefore their linear combination on the right of~\eqref{eq:radical},
satisfies the recurrence. Using $(\sqrt{17})^2=17$, that expression takes
the values $1,8,64$ at $n=0,1,2$. These agree with $a_8(0),a_8(1),a_8(2)$.
Induction in the order-three recurrence proves the identity at every index.
\end{proof}

The first five positive-index values are $8,64,504,3928,30344$, as in
A268943. Formula~\eqref{eq:radical} also shows
\[
 a_8(n)\sim\frac{153+37\sqrt{17}}{34}(3+\sqrt{17})^n,
\]
because $|3-\sqrt{17}|<7<3+\sqrt{17}$ and the displayed leading coefficient
is positive.

\section{Hardin's related table}
Putting $q=k+1$ in~\eqref{eq:rec} proves, uniformly for $k\geq1$,
\begin{equation}\label{eq:column}
 T(n+3,k)=(2k-1)T(n+2,k)+(-k^2+2k+1)T(n+1,k)
                         -k(k+1)T(n,k)\quad(n\geq1).
\end{equation}
These coefficients reproduce all seven empirical column recurrences
listed in A268944, and give the same result for every further column.
For rows one through seven, the corresponding polynomials in $q=k+1$ are
\begin{center}
\begin{tabular}{cl}
\toprule
$n$ & $a_q(n)$ \\
\midrule
1 & $q$ \\
2 & $q^2$ \\
3 & $q^3-q$ \\
4 & $q^4-3q^2+3q$ \\
5 & $q^5-6q^3+11q^2-7q$ \\
6 & $q^6-10q^4+26q^3-30q^2+15q$ \\
7 & $q^7-15q^5+50q^4-81q^3+73q^2-31q$ \\
\bottomrule
\end{tabular}
\end{center}
They follow successively by substitution in~\eqref{eq:rec}, starting with
$1,q,q^2$. Replacing $q$ by $k+1$ gives exactly the seven empirical row
polynomials in A268944. In particular its first three rows are
$k+1,(k+1)^2,k(k+1)(k+2)$, explaining the stated links to A000027,
A000290 and A007531~\cite{rowrefs}.

\begin{corollary}\label{cor:binary}
For the binary column $q=2$,
\[
 S_2(x)=\frac{4x+6x^2}{1-2x^2}-\frac{2x}{1-x}.
\]
Consequently, $a_2(2m)=3\cdot2^m-2$ and
$a_2(2m+1)=2^{m+2}-2$ for $m\geq0$. For $n\geq0$,
\begin{align*}
 a_2(n)&=\sum_{j=0}^n 2^{\min(j,n-j)},\\
 a_2(n+2)&=2a_2(n)+2,\qquad
 a_2(n)-a_2(n-1)=2^{\lfloor n/2\rfloor}\quad(n\geq1).
\end{align*}
These are the corresponding formulas of A027383~\cite{binary}.
The binary words counted here are precisely those for which every substring
has the difference between its numbers of ones and zeros between $-2$ and $2$.
\end{corollary}
\begin{proof}
Substitute $q=2$ in~\eqref{eq:gf} and expand the geometric series to obtain
the parity formulas. Splitting the finite sum at $\lfloor n/2\rfloor$
gives the same parity formulas; subtraction gives the two recurrences.
At $m=0$, the even formula also agrees with our empty-word value $1$.

For the substring interpretation, replace a zero by a step $-1$ and a one
by a step $+1$, and include the initial height zero. Every substring has
the stated bound exactly when the range of all partial sums has width at
most two. Before the first equality event, the steps alternate. That first
event determines an interval of width two: after a $00$ event the path
has descended from its upper endpoint to its lower endpoint, and after a
$11$ event it has ascended in the opposite direction. Between events,
alternating steps keep the path within that interval. The next equality
event stays in the interval exactly when its label is opposite to the
previous event's label. Repeating this argument proves both directions;
with no events the path has width at most one.
\end{proof}

\section{A regular-relation extension of Fried's theorem}
Let $R\subseteq\mathcal A_q\times\mathcal A_q$ be a relation with no
loops, so $(b,b)\notin R$. Suppose that for each $r\in\mathcal A_q$
there are exactly $d$ letters $b$ with $(b,r)\in R$, where
$0\leq d\leq q-1$. Require successive equality-event labels to be related
by $R$, and let $a_R(n)$ count these words, with $a_R(0)=1$.
Put $S_R(x)=\sum_{n\geq1}a_R(n)x^n$.

\begin{corollary}\label{cor:relation}
For every such relation,
\begin{equation}\label{eq:relation}
 S_R(x)=\frac{qx\bigl(1-(q-3)x-(q+d-2)x^2\bigr)}
 {(1-(q-1)x)(1-(q-2)x-(q+d-1)x^2)}.
\end{equation}
In particular, no other features of $R$ affect these counts. At $d=q-1$
this reproves Fried's Theorem~4~\cite{fried}: if $F_q(x)$ denotes its
displayed rational function, then
\[
 1+S_R(x)=F_q(x)=
 \frac{1-(q-3)x-2(q-2)x^2-(q-2)x^3}
 {1-(2q-3)x+(q-1)(q-4)x^2+2(q-1)^2x^3}.
\]
Here the constant term includes the empty word. The positive-index counts
agree with Fried's enumeration. Writing $b_q(n)=a_R(n)$ in this case gives
\begin{equation}\label{eq:friedrec}
 b_q(n+3)=(2q-3)b_q(n+2)+(-q^2+5q-4)b_q(n+1)
             -2(q-1)^2b_q(n)\quad(n\geq1),
\end{equation}
with $b_q(1)=q$, $b_q(2)=q^2$, $b_q(3)=q^3-q$.
This proves the seven empirical column recurrences of A269467~\cite{unequal}.
For $q=k+1$, its rows $n=1,2,3$ are the same as those of A268944;
the next four are
\begin{align*}
 b_{k+1}(4)&=k^4+4k^3+4k^2+k,\\
 b_{k+1}(5)&=k^5+5k^4+7k^3+2k^2-k,\\
 b_{k+1}(6)&=k^6+6k^5+11k^4+4k^3-k^2+k,\\
 b_{k+1}(7)&=k^7+7k^6+16k^5+8k^4-k^3-k.
\end{align*}
Thus all seven displayed row formulas in A269467 follow as well.
\end{corollary}
\begin{proof}
Use $U$ as before, put an eventful word in $A$ when its first letter $b$
and first event label $r$ satisfy $(b,r)\in R$, and put it in $B$ otherwise.
For an eventful suffix with first event $r$, exactly $d$ prepended letters
satisfy $(c,r)\in R$. When the suffix lies in $A$, one of these letters
is its first letter, so there are $d-1$ unequal choices into $A$;
there are $q-d$ unequal choices into $B$, and the equal choice also
enters $B$ since $R$ has no loops. From $B$, there are $d$ choices into
$A$ and $q-d-1$ into $B$, and the equal choice is forbidden. Thus
\[
 U_{n+1}=(q-1)U_n,\quad A_{n+1}=(d-1)A_n+dB_n,\quad
 B_{n+1}=U_n+(q-d+1)A_n+(q-d-1)B_n.
\]
When $d=0$, state $A$ is empty, so its displayed equation still holds.
Writing $D(x)=1-(q-2)x-(q+d-1)x^2$ and solving these two eventful-state
equations gives $A(x)+B(x)=x(1+x)U(x)/D(x)$, which proves
\eqref{eq:relation}. At $d=q-1$, the only possible relation is inequality;
adding one to~\eqref{eq:relation} and expanding gives the displayed
formula of Fried. Comparing coefficients at degree $n+3\geq4$ gives
\eqref{eq:friedrec}; successive substitution gives the displayed rows.
At $d=1$ it gives Theorem~\ref{thm:general}; this also
covers any fixed-point-free permutation of event labels.
\end{proof}

\begin{corollary}\label{cor:events}
Let $c_R(n,r)$ count these words of length $n$ with exactly $r$ equality
events. For $r\geq1$,
\begin{equation}\label{eq:events}
 \sum_{n\geq1}c_R(n,r)x^n=
 \frac{q\,d^{r-1}x^{2r}}
 {(1+x)^{r-1}(1-(q-1)x)^{r+1}},
\end{equation}
where $d^0=1$, including when $d=0$. In particular there are no words
with $n<2r$, and $c_R(2r,r)=q\,d^{r-1}$.
\end{corollary}
\begin{proof}
Mark each equality event by $y$. In the preceding transition, the terms
$U_n$ and the equal-letter contribution $A_n$ entering $B$ acquire a
factor $y$. Solving gives the eventful-word series
\[
 \frac{y x(1+x)U(x)}{(1+x)(1-(q-1)x)-dyx^2}.
\]
Expand its denominator geometrically in $y$ and insert
$U(x)=qx/(1-(q-1)x)$ to obtain~\eqref{eq:events}.
The coefficient at $x^{2r}$ is as stated. Equivalently, a shortest word
with $r$ events is a concatenation of $r$ double letters whose labels
form an $R$-chain; counting backwards gives $q d^{r-1}$ choices.
\end{proof}

\section{Coverage and verification}
\begin{center}
\small
\begin{tabular}{>{\raggedright\arraybackslash}p{0.52\linewidth}>{\raggedright\arraybackslash}p{0.39\linewidth}}
\toprule
Source item & Mathematical coverage \\
\midrule
A268943, FORMULA: Barker's generating function and radical formula
& Proved in Corollary~\ref{cor:barker}. \\
A268943, FORMULA: empirical recurrence
& Proved in Corollary~\ref{cor:barker}. \\
A268944, FORMULA: columns $k=1,\ldots,7$
& Proved for all $k\geq1$ in~\eqref{eq:column}. \\
A268944, FORMULA: rows $n=1,\ldots,7$
& Proved by~\eqref{eq:rec} and the displayed polynomial table. \\
A268944, CROSSREFS: first three rows
& Proved identities with A000027, A000290 and A007531 above. \\
A268944, CROSSREFS: first column; A027383, NAME and elementary FORMULA identities
& Reproved in Corollary~\ref{cor:binary}; prior word proof Fried (2024), Theorem~10, and (2025), Theorem~4. \\
A027383, balanced-substring interpretation
& Proved equivalent to the binary event rule in Corollary~\ref{cor:binary}. \\
Fried's unequal-label enumeration; A269467 column recurrences and seven row formulas
& Reproved; prior proof Fried (2024), Theorem~10, and (2025), Theorem~4. Set $d=q-1$ in Corollary~\ref{cor:relation}. \\
Loopless constant-indegree relations; counts by event number
& Not posed in the two primary entries; proved in Corollaries~\ref{cor:relation}--\ref{cor:events}. \\
A027383's tree, necklace, folding and arithmetic-sequence interpretations; unrelated cross-references
& Not covered. \\
\bottomrule
\end{tabular}
\end{center}

In the \href{https://github.com/the-omega-institute/trureturing/blob/8a9173761a0a8ce40a4238793b5c18e452f76446/D5/S1/Words/Patterns/BarkerCyclicRepetitionCounting.lean}{pinned Lean source},
the declaration \texttt{result} in namespace
\texttt{D5.S1.Words.Patterns.}\allowbreak\texttt{BarkerCyclicRepetitionCounting}
proves \texttt{claim}: the literal eight-letter word count has a unique
zero-constant-term series with the stated generating function, and its
positive-index counts equal~\eqref{eq:radical}. The general-$q$ theorem,
the row polynomials, the binary deductions and the regular-relation and
event-number formulas have the written proofs above. In that same file,
\texttt{count\_recurrence} and \texttt{radical\_count} supply the positive-index
eight-letter recurrence and closed-form identities; their alphabet is fixed at
eight. The empty-word extensions in this note follow from the written proof.
The accompanying Python checker gives finite word enumeration, an independent
forward automaton, exact arithmetic in $\mathbb Z[\sqrt{17}]$ and polynomial
checks; their ranges are specified in \texttt{checks.json}.
The companion \texttt{verify\_siblings.py} checks regular relations,
event-number coefficients, balanced binary substrings and the unequal-label
identities, with ranges in \texttt{sibling-checks.json}.

\begin{thebibliography}{9}
\bibitem{hardin8}
R. H. Hardin, \emph{OEIS A268943}, introduced 16 February 2016;
formulas conjectured by Colin Barker, 17 January 2019,
\href{https://oeis.org/history/view?seq=A268943&v=8}{revision 8},
NAME, FORMULA and EXAMPLE. \url{https://oeis.org/A268943}.
\bibitem{hardintable}
R. H. Hardin, \emph{OEIS A268944}, 16 February 2016,
\href{https://oeis.org/history/view?seq=A268944&v=4}{revision 4},
NAME, FORMULA and CROSSREFS. \url{https://oeis.org/A268944}.
\bibitem{binary}
R. K. Guy and contributors, \emph{OEIS A027383}, revision 193,
20 June 2026, NAME, COMMENTS and FORMULA. \url{https://oeis.org/A027383}.
\bibitem{fried}
Sela Fried, \emph{Proofs of Several Conjectures From the OEIS},
Journal of Integer Sequences \textbf{28} (2025), Article 25.4.3,
Section 3.2, Theorem 4, pp.~4--6.
\url{https://cs.uwaterloo.ca/journals/JIS/VOL28/Fried/fried15.pdf}.
\bibitem{friedpre}
Sela Fried, \emph{Proofs of some Conjectures from the OEIS},
arXiv:2410.07237v3, 26 October 2024, Section~2.8, Theorem~10,
equation~(18), pp.~14--15. \url{https://arxiv.org/abs/2410.07237v3}.
\bibitem{rowrefs}
The OEIS Foundation, \emph{OEIS A000027, A000290 and A007531},
NAME and FORMULA, accessed 14 November 2026.
\url{https://oeis.org/A000027}, \url{https://oeis.org/A000290},
\url{https://oeis.org/A007531}.
\bibitem{unequal}
R. H. Hardin, \emph{OEIS A269467}, FORMULA,
\url{https://oeis.org/A269467}; proved by Fried~\cite[Corollary~5]{fried}.
\end{thebibliography}
\medskip\noindent
\small \copyright\ 2026 Haobo Ma and Wenlin Zhang.
Licensed under \href{https://creativecommons.org/licenses/by/4.0/}{Creative Commons Attribution 4.0 International (CC BY 4.0)}.
\end{document}
